\subsubsection{Server}

El servidor del juego es el componente con autoridad sobre la lógica del mundo, este procesa las acciones de los jugadores, mantiene el estado de
las entidades y distribuye las actualizaciones a los clientes conectados. Está implementado como un proyecto Unity compilado en modo
\textit{headless} (sin interfaz gráfica) y se ejecuta dentro de contenedores Docker desplegados dinámicamente por el \textit{Core Service}. Cada
instancia es responsable de una única zona del mundo, por lo que un mundo con múltiples zonas implica que se ejecuten varias instancias de servidor en paralelo, 
cada una gestionando su propia zona y los jugadores que se encuentren en ella.

\paragraph{Inicializacion}
Cada instancia de zona del servidor se ejecuta dentro de un contenedor Docker. La imagen se construye a partir de una base de Debian mínima
sobre la cual se instalan las dependencias necesarias para ejecutar el binario del servidor de Unity en modo \textit{headless},
junto con las credenciales requeridas (inyectadas vía variable de entorno) para obtener recursos desde el \textit{Core Service} en el momento del arranque.
La imagen expone dos puertos, uno es el 7777/udp para el tráfico del juego y otro es el 7778/tcp para el
\textit{healthcheck}. Al iniciar el contenedor, el servidor recibe como argumentos el identificador del mundo (\texttt{world-id}) y de la zona
 (\texttt{zone-id}) que debe cargar.

\vspace{0.25cm}

Con esos parámetros, el servidor inicia una secuencia de inicialización mediante varios \textit{Loaders}. 
Primero consulta al \textit{Core Service} la información de la zona, por ejemplo, los elementos colocables 
(\textit{Placeables}) y los no colocables (\textit{Creatables}), y la utiliza para cargar los registros y crear los colisionadores.

\vspace{0.25cm}

Una vez completada la carga, el servidor establece conexión con la base de datos no relacional encargada de persistir 
la información del jugador dentro del mundo (última posición en la zona, inventario, cantidad de oro, entre otros datos).

Si dicha base de datos exclusiva para este mundo no existe previamente, se crea dentro del mismo clúster de Mongo
con el correspondiente identificador de mundo.

Luego, inicializa el servicio de \textit{healthcheck}, utilizado por los \textit{Nomad Clients} para verificar el 
estado operativo de la instancia cuyo funcionamiento e integración se describen con mayor detalle en la sección de Infraestructura.

\vspace{0.25cm}

Finalizada esta etapa, el servidor notifica al \textit{Core Service} que la zona se encuentra activa y lista para recibir jugadores.
Si ocurre un fallo en cualquiera de las fases del proceso, la inicialización se cancela y el servidor informa al \textit{Core Service}
el estado fallido de la instancia.

De forma análoga, durante el servidor notifica al \textit{Core Service} que la zona ha 
pasado al estado \textit{offline}.

\paragraph{Game Tick \& Game Loop}

Otro elemento primordial en la capa de \textit{Core} es la implementación del sistema del \textit{GameLoop}.
Este sistema se encarga de manejar el ciclo de actualización del juego, procesamiento de
comandos y encolado de eventos.

\vspace{0.25cm}

Si bien Unity provee un ciclo de actualización mediante funciones como \texttt{Update()} y \texttt{FixedUpdate()}, estas no ofrecen control
suficiente sobre el orden y la sincronización de los sistemas del juego, y generan \textit{overhead} innecesario.
En su lugar, se desactiva la física automática de Unity y se delega
su ejecución al \textit{Game Loop}, que la invoca manualmente en cada \textit{Physics Update} junto con el resto de la lógica del servidor.

Para esto se desarrolló un \textit{Tick Driver}, que se encarga de ejecutar el ciclo principal del servidor y realizar llamadas
al servicio de networking o \texttt{NetworkService} (encargado del envío de eventos) a razón de $60$ ticks por segundo (tps), y al servicio de \textit{game loop}
a razón de $30 tps$. Esto es debido a que la interacción por medio de código con los objetos manejados por Unity, solamente 
se puede realizar desde el hilo principal del proceso (\textit{Main thread}), limitando la posibilidad de realizar un programa
con recepciones y envíos de red en múltiples hilos.

\vspace{0.25cm}

Todos los sistemas que requieren actualizarse en sincronía con el ciclo de juego, como el movimiento, el combate, los \textit{spawners}, entre otros,
se suscriben a un evento que el \textit{Game Loop} emite en cada tick para realizar las actualizaciones.

\paragraph{Limitación de comandos}

Para evitar que el procesamiento de comandos entrantes consuma demasiado tiempo y afecte al resto del ciclo de juego, se aplica un mecanismo
de limitación por tick. 

\vspace{0.25cm}

En cada tick, el servidor procesa como máximo 100 comandos y destina no más de 10 milisegundos a esta tarea. Si se
alcanza cualquiera de estos dos límites antes de vaciar la cola, los comandos restantes no se descartan sino que quedan en espera y se
procesan en el tick siguiente. De esta forma, se garantiza que picos repentinos de actividad, como muchos jugadores enviando acciones al mismo
tiempo, no interrumpan la física ni la lógica del juego, manteniendo el ciclo estable.

\paragraph{Game Systems y State Storage}
Cada entidad del mundo del juego (jugadores, bolsas de botín, NPCs) que requiera syncronizacion con el servidor está representada por un \textit{objeto de red} (\textit{Networked Object}):
una entidad que existe simultáneamente en el servidor y en los clientes de todos los jugadores conectados, manteniéndose sincronizada entre
ellos.

\vspace{0.25cm}

Sobre estos objetos se compone la lógica mediante \textit{Systems}, que son componentes modulares que se adjuntan a la entidad y son
responsables de un aspecto específico de su comportamiento. El \textit{MovementSystem}, por ejemplo, se suscribe al evento de tick del
\textit{Game Loop} y en cada ciclo procesa la dirección de movimiento, aplica la física correspondiente y actualiza la posición del personaje.
Otros sistemas siguen el mismo patrón para manejar el combate, el \textit{stamina}, la salud, entre otros. Esta separación permite que
cada sistema sea independiente y que la lógica del juego se pueda extender sin modificar el resto.

\vspace{0.25cm}

Para que estos sistemas puedan comunicarse y compartir el estado de la entidad, existe un componente dedicado al almacenamiento de estado llamado
\textit{State Storage}. En el caso de los jugadores, el \textit{CharacterStateStorage} centraliza toda la información relevante de la entidad
posición, dirección de movimiento, salud, stamina, ítem equipado, entre otros, y es la única fuente de verdad sobre su estado.

\paragraph{Player Persistence}
La persistencia del jugador es el mecanismo por el cual el estado de un jugador (inventario, oro, posición y progreso de quests) se conserva
en la base de datos y se restaura cada vez que el jugador se conecta a una zona.

\vspace{0.25cm}

Cuando un jugador se une, el servidor consulta su registro en la base de datos y carga sus datos en los sistemas correspondientes. Si el
jugador no tiene datos previos (primera vez que juega), se inicializa con valores por defecto y se crea un nuevo registro.

\vspace{0.25cm}

Durante la sesión, los guardados ocurren en dos modalidades: de forma \textbf{reactiva}, cuando ocurre un cambio relevante (el jugador recoge un ítem,
gasta oro o avanza en una quest o la completa, el cambio se persiste inmediatamente); y de forma \textbf{periódica}, mediante un guardado automático que se ejecuta
cada cierto intervalo de tiempo configurable. Al desconectarse, se realiza un guardado final de todos los datos para garantizar que ningún
progreso reciente se pierda.

